\documentclass[11pt]{article}
\usepackage[T1]{fontenc}
\usepackage{lmodern}
\usepackage{amsmath,amssymb,amsthm}
\usepackage[margin=1in]{geometry}
\usepackage[hidelinks]{hyperref}
\newtheorem{theorem}{Theorem}
\newtheorem{lemma}[theorem]{Lemma}
\newtheorem{proposition}[theorem]{Proposition}
\theoremstyle{remark}
\newtheorem{remark}[theorem]{Remark}
\newcommand{\R}{\mathbb R}
\newcommand{\Z}{\mathbb Z}
\newcommand{\E}{\mathbb E}
\newcommand{\KL}{D}
\newcommand{\F}{\mathcal F}
\newcommand{\risk}{\mathcal R^{\mathrm{seq}}_{n,\alpha}}
\newcommand{\norm}[1]{\left\lVert#1\right\rVert}
\title{Sharp minimax rates for spectral density estimation\protect\\
under sequential local differential privacy}
\author{Kaiwen Jin\thanks{\texttt{yui.ui.siki@gmail.com}}}
\date{13 October 2026}

\begin{document}
\maketitle

\begin{abstract}
We determine the minimax integrated squared error for estimating the spectral
density of a centered stationary Gaussian time series when each observation
is released under local differential privacy and the release may depend on
all earlier private outputs. On a variance-bounded class whose normalized
covariances have Sobolev smoothness $s>1/2$, the risk is, up to constants,
$1\wedge(n\alpha^4)^{-2s/(2s+1)}$, uniformly for $n\ge2$ and
$0<\alpha\le1$. Thus sequential interaction preserves the sharp
non-interactive rate, without a logarithmic loss. The lower bound follows by
representing a covariance perturbation through independent Gaussian
innovations. An innovation that enters only one past raw observation remains
protected by the same privacy inequality after the entire past output has
been released. Gaussian smoothing and the current private channel then give
a divergence bound of order $n\alpha^4$. The matching upper bound is the
non-interactive estimator of Issartel and Roueff.
\end{abstract}

\section{The problem and the minimax rate}\label{sec:problem}

A spectral density describes the covariance of a stationary time series.
Estimating it from locally private measurements is therefore a basic problem
in inference on dependence: each data holder transmits a private version of
one raw observation, while the statistician sees only the transmitted outputs.
Sequential interaction lets a holder consult previous private outputs before
choosing the next release.

Issartel and Roueff \cite[Theorems 1 and 3]{IR} obtain a sharp, log-free
minimax rate for Gaussian spectral density estimation when the releases are
non-interactive. Their lower bound concerns channels that depend only on
their own raw coordinate, and their Section 3.4 leaves the sharp interactive
rate open. Butucea, Klockmann and Krivobokova \cite[Section 2.4]{BKK}
construct sequential estimators; the global minimax question is stated on
p.\,5 of their paper. The difference between these models matters for a
lower bound. With interaction, the conditional law of the current raw
observation depends on the entire private history, whose distribution changes
under the competing spectral densities.

We show that this additional dependence does not improve the joint minimax
rate on the class below. The argument controls a Gaussian innovation through
the full private history, rather than assuming that the privatized coordinates
are conditionally independent. It applies when the privacy parameter is
allowed to shrink with the sample size. The access rule remains essential:
each release may use its own raw observation and previous private outputs,
but cannot read earlier raw observations.

Fix known parameters $s>1/2$, $C_0>1$ and $C_1>0$. For an even, nonnegative,
$2\pi$-periodic spectral density $f$, put
\begin{equation}\label{eq:covariance}
 \gamma_f(h)=\int_{-\pi}^{\pi}f(\lambda)e^{ih\lambda}\,d\lambda,
 \qquad \rho_f(h)=\frac{\gamma_f(h)}{\gamma_f(0)},\qquad h\in\Z.
\end{equation}
The parameter class is
\begin{equation}\label{eq:class}
 \F(C_0,C_1,s)=\left\{f:\ C_0^{-1}\le\gamma_f(0)\le C_0,
 \quad \sum_{h\ne0}|h|^{2s}|\rho_f(h)|^2\le C_1^2\right\}.
\end{equation}
Here and below $f$ is integrable, real and even, and functions are identified
almost everywhere. We use Lebesgue measure for the $L^2([-\pi,\pi])$ norm.
The class and normalization in \eqref{eq:covariance}--\eqref{eq:class} are
exactly those in \cite[(3)]{IR}.

The raw data $X_1,\ldots,X_n$ are consecutive observations of a centered,
real stationary Gaussian process with spectral density $f$. A sequential
mechanism has standard Borel output spaces $E_i$ and Markov kernels
$Q_i(dz_i\mid x_i,z_{<i})$, where $z_{<i}=(z_1,\ldots,z_{i-1})$.
For every history, raw input pair and measurable output event, require
\begin{equation}\label{eq:privacy}
 Q_i(B\mid x,z_{<i})\le e^\alpha Q_i(B\mid x',z_{<i}).
\end{equation}
For a fixed raw vector, its transcript law is the chronological iterated
product of these kernels. Write $\mathcal Q^{\mathrm{seq}}_{n,\alpha}$ for
this class, with the own-observation access rule specified above. Define
\begin{equation}\label{eq:risk}
 \risk=\inf_{Q\in\mathcal Q^{\mathrm{seq}}_{n,\alpha}}
 \inf_{\widehat f}\sup_{f\in\F(C_0,C_1,s)}
 \E_{f,Q}\norm{\widehat f(Z_{1:n})-f}_2^2,
\end{equation}
where the second infimum is over measurable $L^2([-\pi,\pi])$-valued
estimators. The mechanism and estimator may use the known class parameters.

\begin{theorem}\label{thm:main}
For every fixed $s>1/2$, $C_0>1$ and $C_1>0$, there are constants
$0<c\le C<\infty$, depending only on these three parameters, such that
for every integer $n\ge2$ and every $0<\alpha\le1$,
\begin{equation}\label{eq:rate}
 c\left[1\wedge(n\alpha^4)^{-2s/(2s+1)}\right]
 \le \risk\le
 C\left[1\wedge(n\alpha^4)^{-2s/(2s+1)}\right].
\end{equation}
Consequently, along any sequence of admissible parameter pairs, the minimax
risk tends to zero if and only if $n\alpha^4$ tends to infinity.
\end{theorem}

The upper bound is inherited from \cite[Theorem 3]{IR}. The new argument is
a lower bound for every mechanism in \eqref{eq:risk}, including arbitrary
dependence on the private history. On a moving-average pair, the innovation
shared by the current and a past raw observation occurs only once in the
past. Privacy bounds its posterior density relative to its original Gaussian
law by $e^{\pm\alpha}$. Adding the fresh innovation makes the two current
predictive laws close, and privatizing the current coordinate contracts their
divergence again. Each of these two effects contributes a factor of order
$\alpha^2$. A common Gaussian background permits the same calculation for
the Fourier alternatives used in the minimax lower bound.

The proof concerns the precise Gaussian class and access rule above.
It makes no claim for releases that can read past raw data or for
non-Gaussian time series. Section~\ref{sec:innovation} proves the privacy and
smoothing facts, Section~\ref{sec:testing} gives the sequential divergence
estimate, and Section~\ref{sec:minimax} proves Theorem~\ref{thm:main}.

\begin{remark}\label{rem:largen}
Large-sample upper bounds with an $\alpha$-dependent threshold do not by
themselves give a bound uniform in a shrinking privacy parameter.
For example, the clipping calculation in \cite[Appendix A.3, p.\,27]{BKK}
contains the remainder
$\exp[-\widetilde\tau_n\alpha/(4\tau_n^2)]$ with
$\widetilde\tau_n=16(\log n)^{1+\eta}\tau_n^2$, for fixed $\eta>0$.
It equals $\exp[-4\alpha(\log n)^{1+\eta}]$; its polynomial decay therefore
requires an $\alpha$-dependent large-$n$ condition. This qualification is
relevant when comparing their Theorem 7 with \eqref{eq:rate}.
\end{remark}

\section{Privacy of a Gaussian innovation}\label{sec:innovation}

We first record how a single raw coordinate can affect the entire output
history. This observation uses the restriction on access to past raw data.

\begin{lemma}\label{lem:prefix}
Let $L_t(\cdot\mid x_{1:t})$ be the law of the first $t$ outputs of a
mechanism satisfying \eqref{eq:privacy}, with a deterministic raw prefix.
If $x_{1:t}$ and $x'_{1:t}$ differ in only one coordinate, then
\[
 L_t(B\mid x_{1:t})\le e^\alpha L_t(B\mid x'_{1:t})
\]
for every measurable prefix event $B$.
If a latent variable $U$ is independent of a nuisance vector $V$ and the
raw prefix $x(U,V)$ depends on $U$ in at most one coordinate, the channel
from $U$ to the whole private prefix is also $\alpha$-LDP.
\end{lemma}

\begin{proof}
Suppose the changed coordinate is $j$. The histories before $j$ have the
same law. At time $j$, integrate \eqref{eq:privacy} against this common
history law to dominate the joint output measure through time $j$ by
$e^\alpha$ times the other measure. At each later time $k$, both measures
are extended by the same kernel
$z_{<k}\mapsto Q_k(\cdot\mid x_k,z_{<k})$, because $x_k=x'_k$.
The integral of any nonnegative measurable function preserves a measure
domination, so induction gives the claim through time $t$.
For the latent-variable assertion, first fix $V=v$ and apply this comparison
to any two values of $U$. Integrate over the same law of $V$, which is
independent of $U$. If the coordinate never appears, the channel is constant.
\end{proof}

\begin{lemma}\label{lem:posterior}
Let $K$ be an $\alpha$-LDP kernel between standard Borel spaces, with input
law $\pi$ and output law $M=K\pi$. For $M$-almost every output $z$, the
posterior input law has a density $r(\cdot,z)$ relative to $\pi$ satisfying
\begin{equation}\label{eq:posterior}
 e^{-\alpha}\le r(u,z)\le e^\alpha\quad\text{for $\pi$-almost every $u$},
 \qquad \int(r(u,z)-1)^2\pi(du)\le(e^\alpha-1)^2.
\end{equation}
\end{lemma}

\begin{proof}
Averaging the privacy inequality over a second input gives, for every $u$,
$e^{-\alpha}M\le K(\cdot\mid u)\le e^\alpha M$ as inequalities of
measures. Applying these inequalities to the measurable sections of a
product event and integrating in $u$ yields
\[
 e^{-\alpha}(\pi\otimes M)\le \pi\otimes K
       \le e^\alpha(\pi\otimes M).
\]
The Radon--Nikodym derivative $r(u,z)$ of the middle measure has the same
bounds. By Fubini and equality of the output marginals,
$\int r(u,z)\pi(du)=1$ for $M$-almost every $z$; the measures
$r(u,z)\pi(du)$ give a version of the posterior. These uses of
Radon--Nikodym and product integration are covered, for example, by
\cite[Theorems 1.6.6 and 1.6.12]{VH}. Finally
$|r-1|\le e^\alpha-1$ proves the squared-density bound.
\end{proof}

We write $\KL(P\|R)$ for relative entropy and
$\chi^2(P\|R)=\int(dP/dR-1)^2dR$, with value $+\infty$ if the required
absolute continuity fails. The next bound measures the effect of adding an
independent Gaussian innovation to the posterior in
Lemma~\ref{lem:posterior}.

\begin{lemma}\label{lem:smoothing}
Let $G=N(0,1)$, and let $U$ have density $r$ relative to $G$, where
$r-1\in L^2(G)$. Let $W\sim G$ be independent of $U$. If $|t|<1$ and
$P_t$ is the law of $tU+\sqrt{1-t^2}W$, then
\begin{equation}\label{eq:smoothing}
 \chi^2(P_t\|G)\le\frac{t^2}{1-t^2}\int(r-1)^2dG.
\end{equation}
If $e^{-\alpha}\le r\le e^\alpha$ almost everywhere, the same bounds
hold for $dP_t/dG$.
\end{lemma}

\begin{proof}
Under the reference law $U\sim G$, the pair $(U,Y)$, where
$Y=tU+\sqrt{1-t^2}W$, is a bivariate standard normal with covariance
\[
 \Sigma=\begin{pmatrix}1&t\\t&1\end{pmatrix}.
\]
Let $S_t(u,y)$ be its density relative to $G\otimes G$. Its two marginals
are $G$, so the density $p_t=dP_t/dG$ satisfies
\[
 p_t(y)-1=\int(r(u)-1)(S_t(u,y)-1)G(du).
\]
Cauchy--Schwarz and Tonelli give
\[
 \int(p_t-1)^2dG
 \le\left[\int(r-1)^2dG\right]
       \iint(S_t-1)^2\,dG\,dG.
\]
For $x=(u,y)^{\mathsf T}$ the normal density formula gives
\[
 S_t(x)=(\det\Sigma)^{-1/2}
       \exp\!\left[-\tfrac12 x^{\mathsf T}(\Sigma^{-1}-I)x\right].
\]
The matrix $2\Sigma^{-1}-I$ has positive eigenvalues
$(1-t)/(1+t)$ and $(1+t)/(1-t)$, hence determinant one. The change of
variables $(u,y)\mapsto((u+y)/\sqrt2,(u-y)/\sqrt2)$ evaluates the
convergent Gaussian integral as
\[
 \iint S_t^2\,dG\,dG
 =\frac{1}{\det\Sigma\sqrt{\det(2\Sigma^{-1}-I)}}
 =\frac{1}{1-t^2}.
\]
Since $\iint S_t\,dG\,dG=1$, the integral of $(S_t-1)^2$ equals
$t^2/(1-t^2)$. This proves \eqref{eq:smoothing}, including $t=0$.
Finally, $p_t(y)=\int r(u)S_t(u,y)G(du)$ is an average of $r$ under
the reference conditional normal law, so it retains its upper and lower
bounds.
\end{proof}

\section{Sequential testing for a covariance perturbation}\label{sec:testing}

The testing pair below changes one autocovariance and leaves the marginal
variance fixed. A moving-average representation makes the relevant past
innovation fall within Lemma~\ref{lem:prefix}.

\begin{proposition}\label{prop:purelag}
For $a>0$, an integer $h\ge1$ and $|\delta|\le1$, define
\[
 g_\pm(\lambda)=\frac{a}{2\pi}
                 \left[1\pm\frac{\delta}{2}\cos(h\lambda)\right].
\]
Let $M_\pm$ be the output laws of any sequential mechanism satisfying
\eqref{eq:privacy} under these two centered Gaussian observation laws.
For every $n\ge1$ and $0<\alpha\le1$,
\begin{equation}\label{eq:purelag}
 \KL(M_+\|M_-)\le128(n-h)_+\alpha^4\delta^2,
 \qquad (n-h)_+=\max(n-h,0).
\end{equation}
\end{proposition}

\begin{proof}
Choose $A>0$ and signed $B$ by
\begin{equation}\label{eq:ma-coefficients}
 B^2=\frac{1-\sqrt{1-\delta^2/4}}2,
 \qquad A=\sqrt{1-B^2},\qquad B=\frac{\delta}{4A}.
\end{equation}
These relations are consistent and give
$A^2+B^2=1$, $AB=\delta/4$, $B^2\le\delta^2/8$, and $A^2\ge7/8$.
Take independent standard normals $\varepsilon_k$, $1-h\le k\le n$, and
put
\begin{equation}\label{eq:ma}
 U_i^\pm=\sqrt a\,(A\varepsilon_i\pm B\varepsilon_{i-h}),
 \qquad 1\le i\le n.
\end{equation}
The variance is $a$, the lag-$h$ covariance is $\pm aAB=\pm a\delta/4$,
and all other nonzero-lag covariances vanish. These are exactly the Fourier
coefficients of $g_\pm$ in \eqref{eq:covariance}. Centered Gaussian vectors
are determined by their covariances \cite[Section 3.10]{Durrett}, so
\eqref{eq:ma} realizes the observation laws.

Both raw prefix laws are mutually absolutely continuous. Indeed
$g_\pm\ge a/(4\pi)$, and for every real coefficient vector $c$ their
covariance matrices satisfy
\[
 c^{\mathsf T}T_\pm c
 =\int_{-\pi}^{\pi}g_\pm(\lambda)
        \left|\sum_jc_j e^{ij\lambda}\right|^2d\lambda
 \ge\frac a2\sum_j c_j^2.
\]
They therefore have strictly positive Gaussian densities. Passing them
through the same transcript kernel preserves absolute continuity in both
directions. In particular, the two laws of each private history have the
same null sets. All conditional comparisons that follow are made outside
these common null sets; regular conditional laws exist on the standard
Borel spaces involved.

Fix time $i$. The fresh $\varepsilon_i$ is independent of the pair
$(\varepsilon_{i-h},Z_{<i})$: no raw prefix coordinate involves
$\varepsilon_i$, and the prefix kernels use only that prefix. If $i>h$,
the innovation $\varepsilon_{i-h}$ appears in the raw prefix only in
$U_{i-h}^\pm$. Its next appearance is in the current $U_i^\pm$.
Lemma~\ref{lem:prefix}, with all other innovations as independent nuisance
variables, says that its channel to the entire history is $\alpha$-LDP.
Lemma~\ref{lem:posterior} bounds its posterior density relative to $G$ by
$e^{\pm\alpha}$ and its $\chi^2$ divergence by $(e^\alpha-1)^2$,
separately under each sign model.

For a fixed common history $z$, let $P_\pm^z$ be the predictive law of
$U_i^\pm$. Write $G_a=N(0,a)$. Apply Lemma~\ref{lem:smoothing} with
$t=\pm B$; scaling by $\sqrt a$ cancels from both density ratios. Since
$B^2/(1-B^2)\le\delta^2/(8-\delta^2)\le\delta^2/7$, we obtain
\begin{equation}\label{eq:predictive}
 e^{-\alpha}\le\frac{dP_\pm^z}{dG_a}\le e^\alpha,
 \qquad \chi^2(P_\pm^z\|G_a)
        \le\frac{\delta^2}{7}(e^\alpha-1)^2.
\end{equation}
When $i\le h$, $\varepsilon_{i-h}$ is absent from the observed prefix.
It remains a standard normal independent of the history and of the fresh
innovation, so in this case $P_+^z=P_-^z=G_a$.

For $i>h$, put $p_\pm=dP_\pm^z/dG_a$. Using $\log u\le u-1$ and
$p_-\ge e^{-\alpha}$ gives
\begin{align}\label{eq:predictive-kl}
 \KL(P_+^z\|P_-^z)
 &\le \chi^2(P_+^z\|P_-^z)
 \le e^\alpha\int(p_+-p_-)^2dG_a\notag\\
 &\le2e^\alpha\{\chi^2(P_+^z\|G_a)+\chi^2(P_-^z\|G_a)\}
 \le\frac{4e^\alpha}{7}\delta^2(e^\alpha-1)^2.
\end{align}
At the same history $z$, the current kernel $Q_i(\cdot\mid\cdot,z)$ is
one common $\alpha$-LDP channel under both models. By
\cite[Lemma 3]{IR}, for any input laws $P,R$ it satisfies
\begin{equation}\label{eq:channel-contraction}
 \KL(Q_iP\|Q_iR)\le\beta_\alpha\KL(P\|R),
 \qquad \beta_\alpha=\frac12(e^\alpha-e^{-\alpha})^2\wedge1.
\end{equation}
For $0<\alpha\le1$, convexity gives $e^\alpha-1\le(e-1)\alpha<2\alpha$,
and integration of $e^t+e^{-t}\le2e<6$ gives
$e^\alpha-e^{-\alpha}\le6\alpha$. Thus
\[
 \beta_\alpha\frac{4e^\alpha}{7}(e^\alpha-1)^2
 \le\frac{864}{7}\alpha^4<128\alpha^4.
\]
The conditional output divergence at each $i>h$ is at most
$128\alpha^4\delta^2$, and it is zero at each $i\le h$.
The kernel chain rule \cite[Lemma 14]{IR} now gives
\[
 \KL(M_+\|M_-)
 =\sum_{i=1}^n\E_+\!\left[
    \KL\bigl(\mathcal L_+(Z_i\mid Z_{<i})\,
             \|\,\mathcal L_-(Z_i\mid Z_{<i})\bigr)\right],
\]
which proves \eqref{eq:purelag}. Finite products of the output spaces are
standard Borel, so the chain rule and conditional versions apply throughout.
\end{proof}

\begin{theorem}\label{thm:pair}
Let $f_+,f_-$ be even, nonnegative, integrable spectral densities. Suppose
that for some $a>0$, integer $h\ge1$ and $|\delta|\le1$,
\begin{equation}\label{eq:pair-condition}
 \min(f_+,f_-)\ge\frac{3a}{4\pi},\qquad
 f_+-f_-=\frac{a\delta}{2\pi}\cos(h\lambda)
 \quad\text{almost everywhere}.
\end{equation}
Every sequential $\alpha$-LDP mechanism with the specified access rule
satisfies
\begin{equation}\label{eq:pair-kl}
 \KL\bigl(\mathcal L_{f_+,Q}(Z_{1:n})\,
          \|\,\mathcal L_{f_-,Q}(Z_{1:n})\bigr)
 \le128(n-h)_+\alpha^4\delta^2
\end{equation}
for $n\ge1$ and $0<\alpha\le1$.
\end{theorem}

\begin{proof}
With $g_\pm$ as in Proposition~\ref{prop:purelag}, define the common
spectral background
\[
 f_V=\frac{f_++f_-}{2}-\frac{a}{2\pi}\ge\frac{a}{4\pi}.
\]
Let $U^\pm$ be the centered Gaussian vectors with spectral densities
$g_\pm$ constructed in \eqref{eq:ma}, and let $M_\pm$ denote the
transcript laws under $f_\pm$ and $Q$.
The background $f_V$ has a positive definite finite Toeplitz covariance matrix,
by the same
quadratic-form calculation used above. A Cholesky factor therefore constructs
a centered Gaussian vector $V$ with that covariance, independent of $U^\pm$.
The vectors $X^\pm=U^\pm+V$ have the original covariance matrices because
$g_\pm+f_V=f_\pm$. Their centered Gaussian laws are the required raw laws.

Condition on $V=v$. The release kernels become
\[
 Q_i^v(dz_i\mid u_i,z_{<i})
       =Q_i(dz_i\mid u_i+v_i,z_{<i}).
\]
For every fixed $v$, these are still sequential $\alpha$-LDP kernels:
the raw shift preserves \eqref{eq:privacy}, and each kernel still reads only
its current $u_i$ and previous private outputs. Proposition~\ref{prop:purelag}
therefore applies uniformly in $v$. The law $\nu$ of $V$ is common to both
models. The common-mixture divergence inequality \cite[Lemma 15]{IR},
or the joint chain rule followed by projection, gives
\[
 \KL(M_+\|M_-)
 \le\int \KL\bigl(\mathcal L_+(Z_{1:n}\mid V=v)\,
                  \|\,\mathcal L_-(Z_{1:n}\mid V=v)\bigr)\nu(dv).
\]
This proves \eqref{eq:pair-kl}. The conditional transcript kernels are
measurable in $v$ by iterated kernel composition.
\end{proof}

\section{Proof of the minimax bounds}\label{sec:minimax}

\begin{proof}[Proof of Theorem~\ref{thm:main}]
For the lower bound, write $b=n\alpha^4$ and choose
\begin{equation}\label{eq:packing}
 m=\left\lceil(\max(1,b))^{1/(2s+1)}\right\rceil,
 \qquad \kappa=\min\left(\frac{C_1}{\sqrt2},\frac1{512}\right),
 \qquad \theta=\kappa m^{-s-1/2}.
\end{equation}
The choice of $m$ covers both small and large effective sample sizes.
The two restrictions on $\kappa$ enforce the class constraint and keep
neighboring output laws close. For $\omega\in\{-1,1\}^m$, set
\begin{equation}\label{eq:hypercube}
 f_\omega(\lambda)=\frac1{2\pi}
             \left[1+2\theta\sum_{h=1}^m\omega_h\cos(h\lambda)\right].
\end{equation}
Since $s>1/2$, $\theta m\le\kappa$ and
$f_\omega\ge(1-2\kappa)/(2\pi)>0$. These densities are even and periodic.
Direct integration gives $\gamma_\omega(0)=1$,
$\gamma_\omega(\pm h)=\theta\omega_h$ for $1\le h\le m$, and zero at
other nonzero lags. Hence the variance is in $[C_0^{-1},C_0]$, and
\[
 \sum_{h\ne0}|h|^{2s}|\rho_\omega(h)|^2
 =2\theta^2\sum_{h=1}^m h^{2s}\le2\kappa^2\le C_1^2.
\]
Thus all $f_\omega$ are in the original class. Their positive spectral
lower bound makes every finite Toeplitz covariance positive definite.
The corresponding Gaussian finite laws are consistent, so the countable
extension theorem \cite[Theorem A.3.1]{Durrett} gives a stationary Gaussian
process: translation invariance follows from its covariance and Gaussian
identification. This also verifies the observation-law realization of the
subfamily.

If $\omega,\omega'$ differ only at coordinate $h$, then
\[
 f_\omega-f_{\omega'}=\frac{2\theta\omega_h}{\pi}\cos(h\lambda).
\]
Apply Theorem~\ref{thm:pair} with $a=1/4$ and $\delta=16\theta\omega_h$.
Here $|\delta|\le1/32$, and
$\min(f_\omega,f_{\omega'})\ge3/(16\pi)$, as required by
\eqref{eq:pair-condition}. For any admissible mechanism, let $\mu_\omega$
be its entire output law under $f_\omega$. Then
\begin{equation}\label{eq:neighbor}
 \KL(\mu_\omega\|\mu_{\omega'})
 \le32768 b\theta^2
 \le32768\kappa^2\le\frac18,
\end{equation}
because $m^{2s+1}\ge\max(1,b)$.

For any measurable $L^2$-valued estimator, define
\[
 T_h=\operatorname{Re}\int_{-\pi}^{\pi}\widehat f(\lambda)
                     \cos(h\lambda)\,d\lambda,
 \qquad
 \widehat\omega_h=\begin{cases}1,&T_h\ge0,\\-1,&T_h<0.\end{cases}
\]
The coefficient functionals are continuous by Cauchy--Schwarz, so these
decisions are measurable. Since $\int f_\omega\cos(h\lambda)d\lambda
=\theta\omega_h$ and the cosine functions are orthogonal with squared
norm $\pi$, their finite orthogonal projection gives
\begin{equation}\label{eq:loss-reduction}
 \norm{\widehat f-f_\omega}_2^2
 \ge\frac1\pi\sum_{h=1}^m(T_h-\theta\omega_h)^2
 \ge\frac{\theta^2}{\pi}d_H(\widehat\omega,\omega).
\end{equation}
For a complex estimator, first take the real part of its error; this cannot
increase its norm. The second inequality includes zero ties, since a wrong
sign costs at least $\theta$ in its coefficient. Here $d_H$ is the number
of differing sign coordinates.

The Assouad--Pinsker inequality \cite[Lemma 7]{IR} applies to any family
of probability laws on a common measurable space. It states
\[
 \inf_T\sup_\omega\E_\omega d_H(T,\omega)
 \ge\frac m2\min_{d_H(\omega,\omega')=1}
       \left[1-\sqrt{\frac{\KL(\mu_\omega\|\mu_{\omega'})}{2}}\right]
 \ge\frac{3m}{8},
\]
where the last bound is \eqref{eq:neighbor}. Applying it to the actual
sequential output family and using \eqref{eq:loss-reduction} yields, for
every $Q$,
\[
 \inf_{\widehat f}\sup_{f\in\F(C_0,C_1,s)}
        \E_{f,Q}\norm{\widehat f-f}_2^2
 \ge\frac{3\kappa^2}{8\pi}m^{-2s}.
\]
The same inequality holds after taking the infimum over $Q$. As
$m\le2(\max(1,b))^{1/(2s+1)}$, the lower bound in \eqref{eq:rate} holds
with $c=3\kappa^2/(8\pi 2^{2s})>0$.

For the upper bound, \cite[Theorem 3]{IR} supplies a non-interactive
mechanism and a measurable estimator on exactly \eqref{eq:class} with
\[
 \sup_{f\in\F(C_0,C_1,s)}\E_{f,Q}\norm{\widehat f-f}_2^2
 \le C_{C_0,C_1,s}
       \left[1\wedge\{n(\alpha^4\wedge1)\}^{-2s/(2s+1)}\right]
\]
for every $n\ge2$ and $\alpha>0$. Their release adds independent Laplace
noise of scale $3/\alpha$ to the binary indicator
$\mathbf 1_{\{|X_i|>C_0\}}\in\{0,1\}$ and the sign statistic
$\operatorname{sgn}(X_i)\in[-1,1]$. The joint $\ell^1$ sensitivity is
at most $1+2=3$, so this is an $\alpha$-LDP
kernel into $\R^2$. A non-interactive kernel is sequentially admissible by
ignoring the private history. For $\alpha\le1$, the displayed expression
is the upper bound in \eqref{eq:rate}.

Finally, the two-sided estimate shows that the risk tends to zero when
$n\alpha^4\to\infty$. If $n\alpha^4$ fails to tend to infinity, it has a
bounded subsequence; the lower bound is bounded away from zero there.
This proves the converse and completes the theorem.
\end{proof}

The effective size $n\alpha^4$ therefore describes both the consistent
regime and its boundary. In particular, every regime with bounded
$n\alpha^4$ prevents uniformly consistent spectral density estimation,
and no logarithmic factor is required in either minimax bound.

\section*{AI-use statement}
OpenAI Codex was used to develop the argument, retrieve and check cited
source statements, and prepare the manuscript. The mathematical content has
not received independent human verification. Independent mathematical review
is requested.

\begin{thebibliography}{9}
\bibitem{BKK}
C. Butucea, K. Klockmann and T. Krivobokova,
\emph{Nonparametric Spectral Density Estimation using Interactive Mechanisms
under Local Differential Privacy}, Journal of Machine Learning Research
\textbf{27}(168) (2026), 1--49.
\url{https://www.jmlr.org/papers/v27/25-0680.html}.

\bibitem{IR}
Y. Issartel and F. Roueff,
\emph{On the privacy cost for dependent Gaussian data: spectral density
estimation under local differential privacy}, arXiv:2608.24847v1,
25 August 2026.
\url{https://arxiv.org/abs/2608.24847v1}.

\bibitem{Durrett}
R. Durrett, \emph{Probability: Theory and Examples}, fifth edition,
version of 11 January 2019.
\url{https://services.math.duke.edu/~rtd/PTE/PTE5_011119.pdf}.

\bibitem{VH}
R. van Handel, \emph{Stochastic Calculus, Filtering, and Stochastic Control},
lecture notes, version of 29 May 2007.
\url{https://web.math.princeton.edu/~rvan/acm217/ACM217.pdf}.
\end{thebibliography}
\end{document}
